
IPV is not a marginal risk in the population premarital education addresses; it is among the largest measured contributors to disease burden in women of reproductive age. The 2023 Global Burden of Disease modelling study estimated that 608 million (95\% UI 518--724 million) women aged 15 and older had ever experienced IPV, contributing 18.5 million (8.74--30.0) DALYs, and ranked IPV as the fourth leading risk factor for DALYs among women aged 15--49, with anxiety disorders and major depressive disorder the leading attributed outcomes \cite{ref1}. A separate WHO/LSHTM review, pooling 366 studies across 161 countries, estimated that 27\% (95\% UI 23--31\%) of ever-partnered women aged 15--49 have experienced physical or sexual IPV in their lifetime, with the exposure already present in 24\% (21--28\%) of women aged 15--19 and 26\% (23--30\%) of women aged 19--24 \cite{ref2} --- before or around this population's typical marriage age. Regional data confirm marriage itself is not protective: Indonesian surveys found 12.53--17\% of women married before age 18 \cite{ref3,ref4}, and among Rohingya women in Cox's Bazar, child marriage was associated with more than twice the odds of justifying husband-perpetrated abuse (AOR 2.71, 95\% CI 1.78--4.11) and with substantially higher odds of past-year IPV (AOR 1.72, 95\% CI 1.13--2.61) \cite{ref5}.

Despite this burden, the existing evidence base for premarital and marriage-relationship education (MRE) is narrow: built almost entirely to detect effects on relationship satisfaction, communication skill and, at most, marital stability --- not decision control, coercion recognition, or safe help-seeking. The largest MRE meta-analyses report effect sizes for relationship quality ($d\approx 0.30$--$0.36$) and communication skill ($d\approx 0.43$--$0.45$) \cite{ref6}; the most current synthesis of a government-funded programme found small effects on relationship quality ($d=0.114$) and skills ($d=0.132$) and non-significant effects on stability, parenting and child well-being \cite{ref7}, and a meta-analysis of relationship aggression found a marginal overall effect ($d\approx 0.11$), not statistically convincing when restricted to RCTs alone ($d\approx 0.05$) \cite{ref8}. None of these analyses coded a decision-control, coercion-recognition, or safe-help-seeking outcome. Tested against a harder outcome, results are mixed, not reassuring: an 8-year randomized follow-up found no divorce-rate benefit, with couples with pre-existing negative communication patterns \emph{more} likely to divorce after the intervention \cite{ref9}; couples with acute baseline risk, including physical aggression, benefited less than lower-risk couples \cite{ref10}.

Two adjacent literatures sharpen, not close, this gap (detailed in \S4.3). Preconception medicine offers a structured checklist for a premarital contact point \cite{ref11}, and premarital genetic screening shows such a point can be delivered at scale but knowledge gained there does not reliably convert into behaviour: premarital sickle-cell-trait screening uptake across Africa is ``suboptimal'' \cite{ref12}, and premarital thalassaemia education raises knowledge without reliably raising uptake \cite{ref13}. Decision quality is this paper's proposed nearest measurement analogue for \tphe{}'s decision-control construct, detailed in \S4.3.

A second gap concerns population fit. The decision-autonomy and coercion-recognition instruments that exist were built and validated primarily in North American and European populations \cite{ref16,ref17,ref18}; exported to Muslim-majority or Southeast/South Asian settings, cross-cultural transfer has required re-validation and has not always succeeded, e.g. a 13-item empowerment scale adapted to Arabic retained sub-optimal model fit even after revision \cite{ref19}, and a review of 59 studies across 22 countries found persistently poor sexual and reproductive health knowledge among Muslim women worldwide, shaped by religious and cultural framing of information access \cite{ref20}. No decision-autonomy or coercion-recognition instrument in this review has been validated in a Thai population; clinical-safety evidence (\S4.3) adds that cultural-competence training for health professionals repeatedly raises provider knowledge or attitudes without a corresponding patient outcome gain. \tphe{} is proposed into this set of gaps: an evidence base measuring the wrong outcomes for its mechanism, outcome instruments not validated for its population, and adjacent cultural-adaptation evidence showing engagement gains do not reliably transfer into outcome gains.
